\backmatter
\renewcommand{\bibname}{参考文献}
\begin{thebibliography}{99}

\item[] \textbf{■ ローカル教材（本書の執筆で参照したもの）}
\bibitem{kubokawa} 久保川達也．『現代数理統計学の基礎』．共立出版，2017．（第2〜7章を参照）
\bibitem{jss:past} 日本統計学会 編．『日本統計学会公式認定 統計検定1級 公式問題集』（2016〜2019年，2021〜2025年の各年版）．
      統計応用（医薬生物学）および統計数理の問題・解答例．
      （2025年の解答は公式略解に導出過程を補ったもの．本書では「2022年 医薬 問2」「2016年 数理 問2」の形で参照した．）

\item[] \textbf{■ ガイドライン・公的資料（Web で確認）}
\bibitem{ich:e9r1} ICH. E9(R1) Addendum on Estimands and Sensitivity Analysis in Clinical Trials to the Guideline on Statistical Principles for Clinical Trials. 2019.
      \url{https://www.fda.gov/media/108698/download}（FDA 版 2021）
\bibitem{ich:e9} ICH. E9 Statistical Principles for Clinical Trials. 1998.
\bibitem{jp:be} 厚生労働省．後発医薬品の生物学的同等性試験ガイドライン（平成9年12月22日 医薬審第487号，最終改正 令和2年3月19日 薬生薬審発0319第1号）．
      \url{https://www.pmda.go.jp/files/000234565.pdf}
\bibitem{fda:ni} U.S. FDA. Non-Inferiority Clinical Trials to Establish Effectiveness: Guidance for Industry. 2016.
\bibitem{fda:multiple} U.S. FDA. Multiple Endpoints in Clinical Trials: Guidance for Industry. 2022.
\bibitem{ema:missing} EMA. Guideline on Missing Data in Confirmatory Clinical Trials. 2010.
\bibitem{web:lec} LEC東京リーガルマインド．統計検定1級合格講座（試験形式の説明）．\url{https://www.lec-jp.com/jssc/class1/}
\bibitem{web:academaid} Academaid．統計検定1級合格へのロードマップ．\url{https://academ-aid.com/qualif/jssc/roadmap-jssc1}
\bibitem{estimandprimer} Kahan, B.~C. et al. The estimands framework: a primer on the ICH E9(R1) addendum. \textit{BMJ}, 384, e076316, 2024.
      \url{https://pmc.ncbi.nlm.nih.gov/articles/PMC10802140/}

\item[] \textbf{■ 原論文・標準的教科書}
\bibitem{km1958} Kaplan, E.~L. and Meier, P. Nonparametric estimation from incomplete observations. \textit{JASA}, 53, 457--481, 1958.
\bibitem{greenwood1926} Greenwood, M. The natural duration of cancer. \textit{Reports on Public Health and Medical Subjects}, 33, 1--26, 1926.
\bibitem{cox1972} Cox, D.~R. Regression models and life-tables (with discussion). \textit{JRSS B}, 34, 187--220, 1972.
\bibitem{mantel1966} Mantel, N. Evaluation of survival data and two new rank order statistics arising in its consideration. \textit{Cancer Chemotherapy Reports}, 50, 163--170, 1966.
\bibitem{mh1959} Mantel, N. and Haenszel, W. Statistical aspects of the analysis of data from retrospective studies of disease. \textit{JNCI}, 22, 719--748, 1959.
\bibitem{armitage1955} Armitage, P. Tests for linear trends in proportions and frequencies. \textit{Biometrics}, 11, 375--386, 1955.
\bibitem{mcnemar1947} McNemar, Q. Note on the sampling error of the difference between correlated proportions or percentages. \textit{Psychometrika}, 12, 153--157, 1947.
\bibitem{finegray1999} Fine, J.~P. and Gray, R.~J. A proportional hazards model for the subdistribution of a competing risk. \textit{JASA}, 94, 496--509, 1999.
\bibitem{schoenfeld1983} Schoenfeld, D.~A. Sample-size formula for the proportional-hazards regression model. \textit{Biometrics}, 39, 499--503, 1983.
\bibitem{landemets1983} Lan, K.~K.~G. and DeMets, D.~L. Discrete sequential boundaries for clinical trials. \textit{Biometrika}, 70, 659--663, 1983.
\bibitem{obf1979} O'Brien, P.~C. and Fleming, T.~R. A multiple testing procedure for clinical trials. \textit{Biometrics}, 35, 549--556, 1979.
\bibitem{holm1979} Holm, S. A simple sequentially rejective multiple test procedure. \textit{Scandinavian Journal of Statistics}, 6, 65--70, 1979.
\bibitem{littlerubin} Little, R.~J.~A. and Rubin, D.~B. \textit{Statistical Analysis with Missing Data}, 3rd ed. Wiley, 2019.
\bibitem{hanley1982} Hanley, J.~A. and McNeil, B.~J. The meaning and use of the area under a receiver operating characteristic (ROC) curve. \textit{Radiology}, 143, 29--36, 1982.
\bibitem{kp2002} Kalbfleisch, J.~D. and Prentice, R.~L. \textit{The Statistical Analysis of Failure Time Data}, 2nd ed. Wiley, 2002.
\bibitem{agresti2013} Agresti, A. \textit{Categorical Data Analysis}, 3rd ed. Wiley, 2013.
\bibitem{fitzmaurice2011} Fitzmaurice, G.~M., Laird, N.~M. and Ware, J.~H. \textit{Applied Longitudinal Analysis}, 2nd ed. Wiley, 2011.
\bibitem{rothman2008} Rothman, K.~J., Greenland, S. and Lash, T.~L. \textit{Modern Epidemiology}, 3rd ed. Lippincott Williams \& Wilkins, 2008.

\item[] \textbf{■ 第2版で追加した手法の原典（書誌情報のみ．本書の式は本文で独立に導出し，数値計算・シミュレーションで確認した）}
\bibitem{dl1986} DerSimonian, R. and Laird, N. Meta-analysis in clinical trials. \textit{Controlled Clinical Trials}, 7, 177--188, 1986.
\bibitem{higgins2002} Higgins, J.~P.~T. and Thompson, S.~G. Quantifying heterogeneity in a meta-analysis. \textit{Statistics in Medicine}, 21, 1539--1558, 2002.
\bibitem{higgins2003} Higgins, J.~P.~T., Thompson, S.~G., Deeks, J.~J. and Altman, D.~G. Measuring inconsistency in meta-analyses. \textit{BMJ}, 327, 557--560, 2003.
\bibitem{higgins2009} Higgins, J.~P.~T., Thompson, S.~G. and Spiegelhalter, D.~J. A re-evaluation of random-effects meta-analysis. \textit{JRSS A}, 172, 137--159, 2009.
\bibitem{rbg1986} Robins, J., Breslow, N. and Greenland, S. Estimators of the Mantel--Haenszel variance consistent in both sparse data and large-strata limiting models. \textit{Biometrics}, 42, 311--323, 1986.
\bibitem{rubin1987} Rubin, D.~B. \textit{Multiple Imputation for Nonresponse in Surveys}. Wiley, 1987.

\item[] \textbf{■ 他者の分析資料（第2版で参照．個人の申告・意見を含む）}
\bibitem{web:others} 本書の作成者による調査記録「他者分析」（2026年10月24日〜11月1日）．統計検定1級の合格体験記・過去問解説ブログ，X（旧 Twitter）の投稿，GitHub 上の学習リポジトリを収集・要約したもの（ブログ約30本，X 投稿，GitHub リポジトリの日次記録）．
      付録Gでは，記録に含まれる受験者の申告（選択した大問，自己採点，合否）と意見を区別して引用し，出題予想・販売教材の宣伝は改訂の根拠にしていない．個々の記事の URL は調査記録に記載されている．
\end{thebibliography}
